\documentclass[12pt]{article}
\usepackage{amsmath, amssymb, booktabs, geometry}
\geometry{margin=1in}
\setlength{\parskip}{1em}
\setlength{\parindent}{0pt}

\title{Practice Problem Set 3}
\author{ECON 101 -- Macroeconomics}
\date{Winter 2026}

\begin{document}
\maketitle

\section*{Part A – Conceptual Questions}

\textbf{1. Economic Growth vs. Business-Cycle Expansion}  
Explain the difference between long-run economic growth and short-run business-cycle expansion.  
Why is sustained growth primarily determined by productivity rather than total labour hours?



\textbf{2. Productivity and Institutions}  
List three key factors that determine labour productivity and discuss how the financial system contributes to each of them.



\textbf{3. Functions of the Financial System}  
Describe three main roles of the financial system in promoting long-run economic growth.  
Provide an example for each.



\textbf{4. Crowding Out Effect}  
Suppose the government runs a budget deficit. Explain how this affects:
\begin{enumerate}
    \item The market for loanable funds
    \item The equilibrium real interest rate
    \item Private investment
\end{enumerate}



\textbf{5. Phases of the Business Cycle}  
Identify and briefly describe the four phases of the business cycle.  
Which economic indicators usually move together with real GDP during expansions and recessions?

\newpage

\section*{Part B – Quantitative Problems}

\textbf{6. Growth Rate Calculations}  
Real GDP in Canada increased from \$1.8 trillion in 2015 to \$2.2 trillion in 2020.
\begin{enumerate}
    \item Compute the average annual growth rate of real GDP.
    \item Using the Rule of 70, estimate how many years it would take for real GDP to double if this rate continues.
\end{enumerate}



\textbf{7. Per Capita Growth Rate}  
Country A’s real GDP grows by 5\% per year, while its population grows by 2\% per year.
\begin{enumerate}
    \item What is the growth rate of real GDP per capita?
    \item Interpret your result in terms of living standards.
\end{enumerate}



\textbf{8. Market for Loanable Funds}  
The demand for loanable funds is \( D = 800 - 50r \) and the supply is \( S = 200 + 100r \), where \(r\) is the real interest rate (in \%).  
\begin{enumerate}
    \item Find the equilibrium interest rate and quantity of funds.
    \item Suppose government borrowing increases by 100 units at every interest rate (a rightward shift in demand).  
    Determine the new equilibrium interest rate and total funds.
\end{enumerate}



\begin{center}
\rule{0.5\textwidth}{0.4pt}\\
\textit{End of Practice Problem Set}
\end{center}

\end{document}
